\section{Differentials, primitives, and quotients}

For a differential graded Hall construction, each charge carries an
entire cochain complex.  Primitivity of a cycle is an equation between
cochains.  Primitivity of its cohomology class only requires that equation
to hold modulo boundaries.  Height truncation respects both constructions
when the differential preserves charge, but it does not identify them.

\begin{example}[A differential that does not descend]\label{differential-height}
Take
\[
 C=\kk[t]\otimes\Lambda(u),\qquad |t|=0,\quad |u|=1,
 \qquad dt=u,\quad du=0,
\]
with both generators primitive and $\varepsilon(t)=\varepsilon(u)=0$.
This is a differential graded Hopf algebra.  Indeed,
$d(t^m)=mt^{m-1}u$ and $d(t^mu)=0$, so $d^2=0$.
The differential preserves the supercommutative relations.
The coproduct is a chain map on generators and hence on products.
Now assign heights $h(t)=2$ and $h(u)=1$ to the underlying Hopf algebra.
The differential lowers height, and
\[
 t\in F^2C,\qquad dt=u\notin F^2C.
\]
Thus the quotient through height one has no induced differential.
If both generators instead have height one, the differential preserves
height.  Yet the ideal $(t)$ remains unstable, since $dt=u\notin(t)$.
Height preservation and differential stability of a further quotient
ideal are separate requirements.
\end{example}

Let $C=\bigoplus_{\beta\in Q_+}C_\beta$ now be a connected,
coaugmented, counital differential graded coalgebra with
$C_0=\kk1$.  Its differential preserves charge, and
$\Delta1=1\otimes1$, $\varepsilon(1)=1$.
Set $C_+=\ker\varepsilon=\bigoplus_{\beta\neq0}C_\beta$ and
$q=1_C-\eta\varepsilon$.
The reduced coproduct and strict primitive complex are
\begin{equation}
 \bar\Delta=(q\otimes q)\Delta|_{C_+}:C_+\to C_+\otimes C_+,
 \qquad P(C)=\ker\bar\Delta.\label{strict-primitives}
\end{equation}
For $x\in C_+$ this says
$\bar\Delta x=\Delta x-1\otimes x-x\otimes1$.
Taking the kernel on all of $C$ would also include the unit.

\begin{proposition}[Primitive truncation and cohomology]
\label{primitive-comparison}
Put
\[
 P_\beta=\ker\left(C_\beta\xrightarrow{\bar\Delta}
 \bigoplus_{\substack{\gamma+\delta=\beta\\\gamma,\delta\neq0}}
 C_\gamma\otimes C_\delta\right).
\]
Strict primitives commute with height restriction and completion:
\[
 P(\tau_{\leq N}C)=\bigoplus_{0<\htop(\beta)\leq N}P_\beta,
 \qquad P(\widehat C)=\prod_{\beta\neq0}P_\beta.
\]
These are identities of complexes, with the product taken degreewise.
The full cohomology $H^*(C)$ is a charge-graded cohomological coalgebra.
For every $p\in\mathbb Z$ and $\beta\neq0$ there is a natural map
\begin{equation}
 \kappa_\beta^p:H^p(P_\beta)
 \longrightarrow\bigl(P(H^*(C))\bigr)^p_\beta.
 \label{cohomological-primitives}
\end{equation}
Its target is the kernel of the map
\begin{equation}
 H^p(C_\beta)\longrightarrow
 \bigoplus_{\substack{\gamma+\delta=\beta,\ \gamma,\delta\neq0\\r+s=p}}
 H^r(C_\gamma)\otimes H^s(C_\delta)
 \label{full-cohomology-target}
\end{equation}
induced by $\bar\Delta$ and the K\"unneth isomorphism.
The map $\kappa_\beta^p$ need not be injective.
If $C$ is also a differential graded bialgebra, $P(C)$ is a differential
graded Lie algebra under the supercommutator, and the maps $\kappa$ form
a morphism of graded Lie algebras.
\end{proposition}

\begin{proof}
The reduced coproduct is a chain map, so its kernel is a subcomplex.
The primitive equation in charge $\beta$ loses no summand in any window
containing $\beta$.  On the completion, it holds if and only if it holds
in each charge.  These observations prove both identities.

Lemma~\ref{sp:kunnethtensor} gives the natural cohomology cross-product
isomorphism, also for unbounded complexes with the ordinary cochain
tensor product. Applying it to the coproduct defines
\eqref{full-cohomology-target}.  A strict primitive cycle maps to a
primitive class.  A boundary within $P_\beta$ maps to a boundary in
$C_\beta$, which defines \eqref{cohomological-primitives}.
Example~\ref{bar-primitive-example} proves failure of injectivity.

If $x,y$ are homogeneous primitives in a bialgebra, multiplication of
their coproducts with the Koszul interchange gives
\[
 \Delta[x,y]=[x,y]\otimes1+1\otimes[x,y],
 \qquad [x,y]=xy-(-1)^{|x||y|}yx.
\]
The two mixed terms cancel.  Associativity gives graded Jacobi, and
the derivation identity makes $d$ a Lie differential.
Passing to cohomology preserves these identities and the comparison.
\end{proof}

Iterated reduced coproducts with $r$ factors vanish on height $n<r$:
each output factor has positive height.  This is finite-height
conilpotence.  On the completion it holds in every fixed height.
A single infinite series need not have a uniformly vanishing reduced
iterate.  The antipode from Proposition~\ref{connected-antipode} satisfies
$S(x)=-x$ for every primitive.  It preserves charge and does not reverse
a positive root into a negative root.

\subsection{An ordered bar calculation}

Let $D$ be an augmented charge-graded differential graded algebra with
$D_0=\kk$.  Put $D_+=\ker\varepsilon$ and let $sD_+=D_+[1]$, with
$|sa|=|a|-1$.  Its ordered reduced bar construction through height $N$ is
\[
 B_{\leq N}(D)=\kk\oplus
 \bigoplus_{\substack{r\geq1,\ \beta_j\neq0\\
                    \sum_j\htop(\beta_j)\leq N}}
 sD_{\beta_1}\otimes\cdots\otimes sD_{\beta_r}.
\]
All words have length at most $N$. Give this space the differential
and deconcatenation coproduct of Section~\ref{sec:olb-ordered-bar}.

\begin{proposition}[Bar restriction]\label{bar-restriction}
The height restriction $B_{\leq N}(D)=\tau_{\leq N}B_-(D)$ is a
differential graded coalgebra in the total-height tensor category.
These constructions commute with further height restriction.
\end{proposition}

\begin{proof}
Proposition~\ref{prop:olb-bar} proves the differential and coalgebra
identities before truncation. Every local differential operation and
every cut preserves total charge. Proposition~\ref{monoidal-restriction}
therefore restricts the identities, and successive restrictions agree.
\end{proof}

\begin{example}[A strict primitive that becomes a boundary]
\label{bar-primitive-example}
Take $D=\kk[x]$ with $|x|=0$, $h(x)=1$, and $d=0$.
Through height two, its bar coalgebra has
\[
 B^{-2}=\kk[x|x],\qquad
 B^{-1}=\kk[x]\oplus\kk[x^2],\qquad B^0=\kk1,
 \qquad b[x|x]=[x^2].
\]
Both words of length one are strict primitives.
The word $[x|x]$ is not primitive, since
$\bar\Delta[x|x]=[x]\otimes[x]\neq0$.
Thus the strict primitive complex has zero differential, and
$[x^2]$ is nonzero in its degree $-1$ cohomology.
Its image in the full bar cohomology is zero.
This proves the noninjectivity assertion in
Proposition~\ref{primitive-comparison}.

For $D=\kk[x]/(x^2)$ the same deconcatenation term survives, but the
bar differential is zero.  Restriction $B_{\leq2}(D)\to B_{\leq1}(D)$
kills $[x|x]$.  Ordinary tensor projection retains
$[x]\otimes[x]$, with coefficient one, in every characteristic.
Total-height tensor restriction kills this term and commutes with
the coproduct.  This is an ordered associative bar coalgebra.
For a word $[a|b]$ in two independent elements, its reduced coproduct is
$[a]\otimes[b]$, so it need not be cocommutative.
An identification with a chiral coalgebra requires a separate comparison
of tensor categories and operations.
\end{example}

\subsection{Primitive ideals and exact enveloping quotients}

The same failure occurs on a differential graded Hopf carrier.
Let $C=T(\kk a\oplus\kk b\oplus\kk w)$ be the free unital
associative algebra with
\[
 (h(a),h(b),h(w))=(1,2,2),\qquad
 (|a|,|b|,|w|)=(-1,-1,-2),
 \qquad da=db=0,\quad dw=b.
\]
Make $a,b$ primitive and define
\[
 \Delta w=w\otimes1+1\otimes w+a\otimes a.
\]
Extend the differential as a derivation and the coproduct as an algebra
map with Koszul interchange.  Coassociativity on $w$ follows because
both iterates have the extra terms
$a\otimes a\otimes1+a\otimes1\otimes a+1\otimes a\otimes a$.
The chain identity holds on $w$ because $b$ is primitive and $da=0$.
It holds on all generators and therefore on the free algebra.
With augmentation killing the three generators, $C$ is a connected
differential graded bialgebra, hence Hopf.
Its height-two component has basis $a^2,b,w$.
Since $a$ is odd, $\bar\Delta(a^2)=0$, while
$\bar\Delta w=a\otimes a\neq0$.
Thus $P(C)_2=\kk a^2\oplus\kk b$, with zero differential.
The class of $b$ is nonzero in $H^{-1}(P(C)_2)$ but equals the boundary
$dw$ in $C$.  This gives noninjectivity of
\eqref{cohomological-primitives} even for a connected differential
graded Hopf algebra.

\begin{theorem}[Differential-stable primitive quotients]
\label{primitive-quotient}
Let $C$ be a connected charge-graded differential graded bialgebra over
$\kk$.  Let $R\subseteq P(C)$ be homogeneous in charge and cohomological
degree.  Let $J=(R)$ be its two-sided associative ideal.
Assume $dJ\subseteq J$.  Then $C/J$ is a connected differential graded
Hopf algebra, and the quotient preserves antipodes.
The condition $dR\subseteq R$ is sufficient for differential stability.
Its completion is the topological differential graded Hopf quotient
\begin{equation}
 \widehat{C/J}\cong\widehat C/\overline J,
 \qquad \overline J=\prod_{\beta\in Q_+}J_\beta.
 \label{completed-primitive-quotient}
\end{equation}
Finite height restrictions commute with this quotient in their
truncated tensor categories.
\end{theorem}

\begin{proof}
For $r\in R$, $\Delta r=r\otimes1+1\otimes r$ and
$\varepsilon(r)=0$.  Multiplicativity gives
\[
 \Delta J\subseteq J\otimes C+C\otimes J,
 \qquad \varepsilon J=0.
\]
These inclusions hold termwise on each product $arb$, with the Koszul
signs.  Thus coproduct and counit descend to the associative quotient.
The assumed differential stability makes it a differential graded
bialgebra.  If $dR\subseteq R$, the Leibniz rule proves $dJ\subseteq J$.
Since $R$ has positive height, $J_0=0$, so the quotient is connected.
Proposition~\ref{connected-antipode} supplies its antipode.
The quotient map intertwines antipodes by uniqueness of convolution
inverses in the convolution algebra of maps $C\to C/J$.

The product of the componentwise quotient maps is surjective, with
kernel $\prod_\beta J_\beta$.  Its target has the quotient height
topology.  This kernel is the closure of $\bigoplus_\beta J_\beta$,
since every finite truncation of a point in the product belongs to $J$.
This proves \eqref{completed-primitive-quotient}, including the
continuous structure maps.  Every finite component of $J$ is generated
using factors of no greater height, so truncation commutes with forming
the indicated quotient ideal.
\end{proof}

\begin{example}[New primitives after quotient]\label{new-primitives}
Consider the ordinary commutative algebra $C=\kk[x,y,z]$, with
$h(x)=h(y)=1$ and $h(z)=2$.  Give $x,y$ primitive coproducts and put
\[
 \Delta z=z\otimes1+1\otimes z+x\otimes y,
 \qquad\varepsilon(x)=\varepsilon(y)=\varepsilon(z)=0.
\]
Extend $\Delta$ as an algebra map.  Both iterated coproducts of $z$
contain, besides its three primitive terms,
\[
 x\otimes y\otimes1+x\otimes1\otimes y+1\otimes x\otimes y.
\]
Thus $C$ is a connected graded bialgebra and hence Hopf.
In the height-two basis $(x^2,xy,y^2,z)$, the four reduced coproducts are
\[
 2x\otimes x,\quad x\otimes y+y\otimes x,\quad
 2y\otimes y,\quad x\otimes y.
\]
They are linearly independent, so there are no height-two primitives.
In the Hopf quotient $C/(x)=\kk[y,z]$, the class of $z$ becomes
primitive.  Therefore the primitive ideal theorem alone does not
identify $P(C/J)$ with a quotient of $P(C)$.
\end{example}

An enveloping presentation supplies the missing comparison.
The characteristic-zero hypothesis is decisive here.

\begin{proposition}[Primitives of an enveloping quotient]
\label{enveloping-quotient}
Let $\mathfrak g$ be a positively charge-graded differential graded Lie
algebra over $\kk$.  Let $\mathfrak r$ be a homogeneous differential
graded Lie ideal.  Give $U(\mathfrak g)$ its primitive coproduct.
There are natural isomorphisms
\begin{align*}
 U(\mathfrak g)/(\mathfrak r)&\cong U(\mathfrak g/\mathfrak r),\\
 P\bigl(U(\mathfrak g)/(\mathfrak r)\bigr)&\cong
 \mathfrak g/\mathfrak r,\\
 P\bigl(\widehat{U(\mathfrak g)/(\mathfrak r)}\bigr)&\cong
 \prod_{\beta\neq0}(\mathfrak g/\mathfrak r)_\beta.
\end{align*}
The first is a differential graded Hopf isomorphism.  The second and
third preserve Lie brackets and differentials, with the height topology
in the third case.
\end{proposition}

\begin{proof}
Corollary~\ref{sp:quotient} constructs the inverse differential graded
Hopf maps in the first line. Theorem~\ref{sp:primitives} identifies the
primitives with $\mathfrak g/\mathfrak r$ using the signed PBW
symmetrization. Proposition~\ref{sp:completion} gives the degreewise
height product in the third line. All identifications preserve the
Lie bracket and differential.
\end{proof}

\subsection{Relation ideals and the choice of coproduct}
\label{sec:olb-relation-ideals}

For a coalgebra $A$, a subspace $J$ is a coideal when
$\Delta J\subseteq J\otimes A+A\otimes J$ and
$\varepsilon(J)=0$.  An associative ideal need not satisfy
these conditions.  A quotient of a bialgebra $A$ by an ideal
$J$ inherits its coproduct and counit precisely when
\[
 \Delta J\subseteq J\otimes A+A\otimes J,
 \qquad \varepsilon(J)=0.
\]
Writing $q:A\to A/J$ for the quotient, these conditions say that
$(q\otimes q)\Delta$ and $\varepsilon$ vanish on $J$. They therefore
factor through $q$ to give the quotient coproduct and counit.  In the differential graded case one
also needs $dJ\subseteq J$ to define the quotient differential.

Take $A=\kk[x]$, with $x$ even and primitive.
The ideal $(x^2)$ is an associative ideal, but
\[
 \Delta(x^2)=x^2\otimes1+2x\otimes x+1\otimes x^2.
\]
Modulo that ideal in both factors, the middle term is nonzero.
Indeed $1,x$ are independent in $\kk[x]/(x^2)$, and $2\ne0$.
Thus the quotient does not inherit the primitive coproduct.

For the odd-square algebra $x^2=t$, the ideal $(t-1)$ fails
even the augmentation condition, since $\varepsilon(t-1)=-1$.
Both quotients exist as associative algebras.

The special quotient in Corollary~\ref{sp:quotient} has a
different hypothesis: its ideal is generated by a differential
graded Lie ideal of primitive elements.
The primitive equations give the coideal condition directly,
and stability of that Lie ideal gives differential stability.
That is why the conclusion is an enveloping Hopf quotient.

This distinction also applies to differential algebras with
nonzero odd squares.  For $r\geq1$, form
\[
 A_r=\kk[t]\langle\tau\rangle[z]/(z^2),\qquad
 |t|=0,\quad|\tau|=-1,\quad|z|=2,
\]
where $t,z$ commute with every element.
Its basis consists of $t^a\tau^bz^c$, with $a,b\geq0$ and
$c=0,1$.  Define
\[
 d\tau=t^r-2\tau^2z,\qquad dt=dz=0.
\]
The right side commutes with $\tau$, so $d(\tau^2)=0$ and
$d^2\tau=0$.

The differential preserves all the relations, giving a unital
differential graded algebra. Its augmentation sends $t,\tau,z$ to
zero. This map commutes with the differential because $r\geq1$.
Making all three displayed generators primitive would not
give a differential graded Hopf presentation.
Already the relation $z^2=0$ would have reduced coproduct
$2z\otimes z\ne0$.
Thus the enveloping cohomology theorem does not replace the
associative higher-operation analysis of this algebra.

There is nevertheless an actual functorial implication for
Lie comparisons.  If $f:\mathfrak g\to\mathfrak h$ is a
quasi-isomorphism of differential graded Lie algebras over
$\kk$, then $U(f)$ is a quasi-isomorphism.
Indeed the natural isomorphisms of Theorem~\ref{sp:cohomology}
identify its cohomology map with $U(H(f))$, which is an
isomorphism because $H(f)$ has an inverse Lie map.
A Lie quasi-isomorphism is therefore sufficient input.
A vector-space identification of cohomology, without that
Lie map, is different data.

